%Version 3.1 December 2024
%%%%%%%%%%%%%%%%%%%%%%%%%%%%%%%%%%%%%%%%%%%%%%%%%%%%%%%%%%%%%%%%%%%%%%
%%                                                                 %%
%% Title Page for Double-Blind Peer Review                        %%
%% European Journal of Psychology of Education (EJPE)             %%
%%                                                                 %%
%%%%%%%%%%%%%%%%%%%%%%%%%%%%%%%%%%%%%%%%%%%%%%%%%%%%%%%%%%%%%%%%%%%%%

\documentclass[pdflatex,sn-apa]{sn-jnl}% APA Reference Style

%%%% Standard Packages
\usepackage{graphicx}%
\usepackage{amsmath}%
\usepackage[utf8]{inputenc}%
\usepackage[T1]{fontenc}%
\usepackage{hyperref}

\begin{document}

\title[Political Attitudes and Historical Reasoning]{Do Political Attitudes Contaminate Historical Reasoning? Ideological Attitudes and Disciplinary Thinking in Czech Adolescents}

\author*[1,2]{\fnm{Juda} \sur{Kaleta}}\email{juda.kaleta@ff.cuni.cz}

\affil*[1]{\orgdiv{Institute of History, Faculty of Arts}, \orgname{Charles University}, \orgaddress{\city{Prague}, \country{Czech Republic}}}

\affil[2]{\orgdiv{Laboratorio Temporal, Faculty of Education}, \orgname{Universidad de Murcia}, \orgaddress{\city{Murcia}, \country{Spain}}}

\abstract{When educational assessments use morally and politically charged historical content, students' personal political attitudes could interfere with their cognitive performance, producing construct-irrelevant variance in test scores. This study tests whether right-authoritarian attitudes predict performance on a widely used historical perspective taking (HPT) instrument whose scenario---a young man in the Weimar Republic considering joining an extremist movement---creates a potential alignment between ideological sympathy and ostensibly ``correct'' disciplinary responses. We administered measures of HPT, right-authoritarian ideology, authoritarianism, historical knowledge, and social desirability to 293 Czech secondary students across 10 schools and 20 classrooms. Two preregistered hypotheses predicted that stronger right-authoritarian attitudes would be associated with higher HPT scores (H1) and that this association would persist after controlling for knowledge and social desirability (H2). Neither hypothesis was supported. Ideology was unrelated to HPT in zero-order correlations ($r = -.05$ to $.01$), multilevel models ($\beta \approx 0$, all $p$s $> .43$), and equivalence tests (all TOST $p$s $< .003$). Multi-group confirmatory factor analysis demonstrated full scalar measurement invariance across ideology groups, and no items exhibited differential functioning. Historical knowledge was the only consistent predictor of HPT performance ($\beta = 0.21$--$0.34$, $p < .001$). These findings indicate that disciplinary historical reasoning, as captured by this instrument, operates independently of students' political attitudes in mainstream classroom settings---a result with implications for educational assessment design in politically sensitive domains.}

\keywords{historical reasoning, political attitudes, construct validity, measurement invariance, multilevel modeling, educational assessment}

\maketitle

%%==================================%%
%% ORCID                           %%
%%==================================%%

\noindent\textbf{ORCID}

\medskip
\noindent Juda Kaleta: \url{https://orcid.org/0000-0003-0043-4039}

%%==================================%%
%% Current Themes of Research and  %%
%% Most Relevant Publications      %%
%%==================================%%

\section*{Current Themes of Research and Most Relevant Publications}

\subsection*{Current Themes of Research}

The author's research focuses on history education and educational assessment. Current work examines the construct validity of disciplinary reasoning assessments in history education, including how political attitudes, task design, and content framing may introduce construct-irrelevant variance. Additional research interests include curriculum analysis and the application of computational methods to educational research.

\subsection*{Most Relevant Publications}

\begin{enumerate}
    \item Kaleta, J. (2025). Student perspectives on high-stakes testing (2015--2025): An AI-assisted systematic review of empirical evidence on motivation, learning, and well-being. \url{https://doi.org/10.71240/lcyc.736516}
    \item Kaleta, J. (2025). Sticking to tradition: History content distribution in Czech school curricula. \textit{History Education Research Journal}. \url{https://doi.org/10.14324/HERJ.22.1.14}
    \item Kaleta, J. (2024). Evaluating the potential of LLMs for thematic analysis in Czech history curricula. \url{https://doi.org/10.31235/osf.io/dtj6h}
\end{enumerate}

\subsection*{Motivation for Submitting to EJPE}

This manuscript addresses a question at the intersection of educational psychology and disciplinary assessment: whether students' political attitudes introduce construct-irrelevant variance into cognitive assessments that use politically sensitive content. The study draws on established frameworks from educational measurement (construct validity, measurement invariance) and social psychology (identity-protective cognition, authoritarianism) to examine a threat to assessment fairness that has implications beyond the specific domain of history education. The European Journal of Psychology of Education, with its focus on the psychological dimensions of learning and assessment across diverse educational contexts, is well suited for this interdisciplinary contribution.

%%==================================%%
%% Declarations                    %%
%%==================================%%

\section*{Statements and Declarations}

\subsection*{Competing Interests}

The author declares no financial or non-financial competing interests.

\subsection*{Funding}

This research received no external funding.

\subsection*{Ethics Approval}

This study was conducted in accordance with the ethical requirements of Charles University, Faculty of Arts, and applicable Czech legal provisions for research involving minors. Research procedures were reviewed and approved by the ethics committee of the Institute of History prior to data collection.

\subsection*{Consent to Participate}

Participation was fully voluntary; written informed consent was obtained from parents or legal guardians of all participants, and each student's assent was confirmed before administration.

\subsection*{Data Availability}

De-identified data and analysis scripts are available on the Open Science Framework (\url{https://osf.io/yng37/}).

\subsection*{Code Availability}

All analysis code is available on the Open Science Framework (\url{https://osf.io/yng37/}).

\subsection*{Author Contribution}

J.K.\ is the sole author and was responsible for all aspects of the study, including conceptualization, methodology, data collection, formal analysis, visualization, and writing.

\subsection*{Preregistration}

Hypotheses, instruments, and analysis plan were preregistered on the Open Science Framework prior to data collection (\url{https://osf.io/yng37/}).

\end{document}
